\documentclass[12pt]{article}
\usepackage[utf8]{inputenc}
\usepackage[french]{babel}
\usepackage{amsmath}
\usepackage{tikz}
\usepackage{circuitikz}
\usepackage{pgfplots}

\begin{document}

\section*{Circuit RC Simplifié}

\subsection*{Schéma}
\begin{center}
\begin{tikzpicture}[american, scale=1]
    % Générateur
    \draw (0,0) node[ground]{} -- ++(0,3) node[ocirc]{} -- ++(0,1) to[V] ++(0,-1) to[short] (0,0);
    % Résistance
    \draw (0,1) to[R] ++(2,0) node[midway,above] {$R$};
    % Condensateur
    \draw (2,0) to[C] ++(0,3) node[midway,above] {$C$} -- (0,3);
    % Annotation
    \node at (1,2) {$V_{in}$};
\end{tikzpicture}
\end{center}

\subsection*{Graphique Théorique}
Voici la courbe théorique de charge avec $RC = 0.005$ s :
\begin{center}
\begin{tikzpicture}
\begin{axis}[
    xlabel={Temps (s)},
    ylabel={Tension (V)},
    xmin=0, xmax=0.005,
    ymin=0, ymax=10,
    samples=200
]
\addplot[blue] {10*(1-exp(-x/0.005))};
\end{axis}
\end{tikzpicture}
\end{center}

\subsection*{Équations}
La tension aux bornes du condensateur :
\[ V_C(t) = V_S \left(1 - e^{-t/RC}\right) \]

\end{document}